
%===================================================
% Информация о типе работы
%===================================================

\newif\ifVkrOfBachelor  %Флаг для определения типа работы
%\VkrOfBachelortrue     %работа бакалавра -  магистрам закоментировать!


%Тена вкр
\def\vkrTheme{Исследование методов автоматического составления расписания и разработка программного средства поддержки генерации расписания для МГТУ «СТАНКИН»}
%Направление
\def\vkrMajor{09.04.04} %09.03.01 09.03.04  09.04.01  09.04.04
\def\vkrMajorName{Программная инженерия}
%Профиль
\def\vkrSpesialization{Технологии разработки интеллектуальных систем и программных комплексов}
%

%Кафедра 
\def\vkrKafedra{информационных технологий и вычислительных систем}
\def\vkrKafedraShort{ИТиВС}

%Институт 
\def\vkrInstitute{информационных технологий}
\def\vkrInstituteShort{ИТ}

%Зав. кафедрой 
\def\vkrKafedraHead{Новоселова Ольга Вячеславовна}
\def\vkrKafedraHeadShort{Новоселова О.\,В.}
\def\vkrKafedraHeadTitles{к.т.н., доц.}

%===================================================
% Информация о студенте
%===================================================
\def\vkrStudentGroup{ИДМ-24-08}
\def\vkrStudent{Кочановский Кирилл Андреевич}
\def\vkrStudentR{Кочановского Кирилла Андреевича} %Родительный падеж
\def\vkrStudentShort{Кочановский К.\,А.}   % \, - короткий пробел 


%===================================================
% Информация о научном руководителе
%===================================================

\def\vkrAdvisor{Гаврилов Андрей Геннадьевич}
\def\vkrAdvisorShort{Гаврилов А.\,Г.}

%должность 
\def\vkrAdvisorPosition{доцент}  %ст. преподаватель, доцент, профессор или зав. кафедрой
\def\vkrAdvisorPositionShort{доц.}  %ст. преп., доц., проф., зав. каф.,

%ученое степень
\def\vkrAdvisorDegree{кандидат технических наук}  %кандидат или доктор  такиих то наук или пусто
\def\vkrAdvisorDegreeShort{к.т.н.}  %аббривииатура первых букв степени или пусто

%ученое звание 
\def\vkrAdvisorTitle{доцент}  %доцент, профессор (не путать с должностью!!, могут совпадать и отличаться) или пусто
\def\vkrAdvisorTitleShort{доц.}  %доц., проф. или пусто


%===================================================
% Информация о научном консультанте (для магистров)
%===================================================

\def\vkrConsultant{Тохунц Наталия Борисовна}
\def\vkrConsultantShort{Тохунц Н.\,Б.}

%должность 
\def\vkrConsultantPosition{доцент}  %ст. преподаватель, доцент, профессор или зав. кафедрой
\def\vkrConsultantPositionShort{доц.}  %ст. преп., доц., проф., зав. каф.,

%ученое степень
\def\vkrConsultantDegree{}  %кандидат или доктор  такиих то наук или пусто
\def\vkrConsultantDegreeShort{}  %аббривииатура первых букв степени или пусто
%ученое звание 
\def\vkrConsultantTitle{} %доцент, профессор (не путать с должностью!!, могут совпадать и отличаться) или пусто
\def\vkrConsultantTitleShort{}  %доц., проф. или пусто




